\title{DeepSeekMath: Pushing the Limits of Mathematical Reasoning in Open Language Models}

\begin{abstract}
The intricate and highly structured demands of mathematical reasoning present notable difficulties for language models. In this study, we present DeepSeekMath~7B, which extends the pre-training of DeepSeek-Coder-Base-v1.5 7B by incorporating 120B mathematics-focused tokens derived from Common Crawl, integrated with natural language and programming code materials. DeepSeekMath~7B attains an outstanding 51.7\% accuracy on the challenging MATH benchmark without the assistance of external tools or ensemble techniques, bringing its performance close to leading models such as Gemini-Ultra and GPT-4. Furthermore, applying self-consistency using 64 samples from DeepSeekMath~7B results in a 60.9\% score on MATH. The remarkable mathematical reasoning performance of DeepSeekMath~arises from two crucial aspects: first, we exploit the extensive capabilities of publicly accessible web resources through a carefully crafted data selection procedure; second, we develop Group Relative Policy Optimization (GRPO), an adaptation of Proximal Policy Optimization (PPO), which not only boosts mathematical reasoning proficiency but also improves PPO's memory efficiency.
\end{abstract}

\section{Introduction}

Large language models (LLMs) have transformed artificial intelligence’s methods for mathematical reasoning, resulting in notable progress on both quantitative reasoning benchmarks \citep{MATH} and geometry reasoning evaluations \citep{trinh2024solving}. In addition, these systems have demonstrated substantial utility in aiding humans to tackle intricate mathematical challenges \citep{tao}. Despite these advances, state-of-the-art models like GPT-4 \citep{gpt4} and Gemini-Ultra \citep{gemini} remain proprietary, while open-source alternatives that are available continue to lag significantly in terms of performance.

This paper presents DeepSeekMath, a specialized language model tailored for mathematical reasoning, which delivers substantially stronger mathematical performance than other open-source models and comes close to matching GPT-4 on established academic tests. Central to this advancement is the construction of the DeepSeekMath Corpus—a vast, high-quality pre-training dataset containing 120B tokens drawn from mathematical content. The corpus was assembled from Common Crawl (CC) via a classifier based on fastText \citep{joulin2016fasttext}. During the initial phase, the classifier was trained with samples from OpenWebMath \citep{openwebmath} as positive examples, while an assortment of unrelated web pages acted as negative examples. Afterward, this classifier was used to extract further positive samples from CC, which underwent refinement with human annotation. These new annotations allowed us to retrain and enhance the classifier. The quality of the resulting dataset is reflected in our evaluation: DeepSeekMath-Base 7B, our foundational model, attained 64.2\% on the GSM8K benchmark \citep{gsm8k} and 36.2\% on the MATH competition dataset \citep{MATH}, surpassing Minerva 540B \citep{minerva}. Notably, the DeepSeekMath Corpus encompasses multiple languages, yielding measurable gains on Chinese mathematical tasks \citep{wei2023cmath,agieval}. We hope our methodology in mathematical data processing serves as a useful reference for the research community and anticipate considerable opportunities for future improvement.

DeepSeekMath-Base is built upon DeepSeek-Coder-Base-v1.5 7B \citep{deepseek-coder}, as preliminary evidence suggests that leveraging a model initially trained on code yields superior results relative to models based on general-purpose LLMs. Additionally, our experiments reveal that mathematical training also leads to improved performance on MMLU \citep{mmlu} and BBH benchmarks \citep{bbh}. This suggests that mathematical fine-tuning not only advances mathematical proficiency but also bolsters general reasoning skills within the model.

Following pre-training, we perform mathematical instruction tuning on DeepSeekMath-Base, utilizing datasets incorporating chain-of-thought \citep{cot}, program-of-thought \citep{pot,pal}, and tool-integrated reasoning \citep{tora} approaches. The finalized model, DeepSeekMath-Instruct 7B, surpasses other models of similar (7B) size and achieves performance on par with instruction-tuned open-source models at the 70B scale.

In addition, we propose Group Relative Policy Optimization (GRPO), a reinforcement learning (RL) algorithm derived from Proximal Policy Optimization (PPO) \citep{schulman2017proximal}. Unlike traditional approaches, GRPO eliminates the need for a critic model by leveraging group scores to estimate the baseline, thereby markedly reducing computational demands during training. Using only a portion of English instruction-tuning data, GRPO demonstrates notable performance gains over the robust DeepSeekMath-Instruct across a variety of benchmarks. Specifically, it achieves significant improvements on in-domain tasks (GSM8K: 82.9\% $\rightarrow$ 88.2\%, MATH: 46.8\% $\rightarrow$ 51.7\%) as well as on out-of-domain mathematical benchmarks (e.g., CMATH: 84.6\% $\rightarrow$ 88.8\%) during reinforcement learning. Furthermore, we establish a unified framework that connects diverse methodologies, including Rejection Sampling Fine-Tuning (RFT) \citep{yuan2023scaling}, Direct Preference Optimization (DPO) \citep{dpo}, PPO, and GRPO, revealing that these algorithms can be framed as forms of either direct or streamlined RL strategies. We conduct comprehensive experiments—such as comparisons between online and offline learning, outcome versus process-based supervision, and single-turn against iterative RL—to systematically analyze the critical components underlying this framework. Finally, we elucidate how our RL approach enhances the performance of instruction-tuned models and outline future research opportunities for developing even more effective RL methods within this consolidated framework.

\subsection{Contributions}
We present a scalable approach to math pre-training and further investigate and analyze the application of reinforcement learning.

\textbf{Math Pre-Training at Scale}

\begin{itemize}[topsep=0pt]
    \item This study demonstrates that Common Crawl’s publicly available dataset encompasses substantial mathematically relevant content. Utilizing a rigorously structured data selection pipeline, we assembled the DeepSeekMath Corpus—a curated dataset consisting of 120B tokens drawn from web sources focused on mathematics. This resource is nearly sevenfold greater in size than the math-specific web dataset employed by Minerva \citep{minerva} and approximately nine times larger than the recently launched OpenWebMath \citep{openwebmath}.

\item DeepSeekMath-Base 7B, our pre-trained base model, demonstrates performance on par with Minerva 540B \citep{minerva}. This suggests that model size alone does not solely determine effectiveness in mathematical reasoning; smaller models trained on superior data can also deliver robust results.

\item We present insights from our experiments on training with mathematical tasks. Pretraining models on code before engaging in math-specific training enhances their performance on mathematical problem-solving, regardless of whether external tools are utilized. This provides some evidence addressing the ongoing debate: \textit{does training on code bolster reasoning skills?} Our results indicate that it does, particularly concerning mathematical reasoning.

\item While it is a widespread practice to use arXiv papers for training, particularly in mathematics-focused research, this approach does not yield significant gains on any of the mathematical evaluation benchmarks utilized in this study.
\end{itemize}

\textbf{Exploration and Analysis of Reinforcement Learning}

\begin{itemize}[topsep=0pt]
    \item We present Group Relative Policy Optimization (GRPO), a reinforcement learning approach that is both efficient and robust. Unlike traditional methods such as Proximal Policy Optimization (PPO) that rely on a critic model, GRPO estimates the baseline using group scores, thereby greatly decreasing the computational and training demands.
    \item Our results show that applying GRPO notably boosts the capabilities of our instruction-tuned model DeepSeekMath-Instruct, utilizing only instruction-tuning data for training. Additionally, improvements in generalization to out-of-domain tasks are observed throughout the reinforcement learning phase.
    \item We propose a cohesive framework that encompasses a range of methodologies including RFT, DPO, PPO, and GRPO. Comprehensive experiments are carried out—such as comparisons between online and offline training, outcome-based versus process-based supervision, and analyses of single-turn versus iterative reinforcement learning—to thoroughly examine the critical factors underpinning this framework.
    \item Building on the unified paradigm, we investigate the mechanisms by which reinforcement learning proves effective and outline several promising avenues for further enhancing reinforcement learning techniques for large language models.
\end{itemize}

\subsection{Summary of Evaluations and Metrics}

\begin{itemize}[topsep=0pt]
    \item \textbf{English and Chinese Mathematical Reasoning}: We perform thorough evaluations of our models using both English and Chinese datasets that span mathematical questions from elementary through collegiate difficulty. For English-language tasks, we incorporate GSM8K \citep{gsm8k}, MATH \citep{MATH}, SAT \citep{llemma}, OCW Courses \citep{minerva}, and MMLU-STEM \citep{mmlu}. The Chinese evaluation suite comprises MGSM-zh \citep{mgsm}, CMATH \citep{wei2023cmath}, Gaokao-MathCloze \citep{agieval}, and Gaokao-MathQA \citep{agieval}. Our assessment measures the proficiency of models in producing independent written solutions as well as solving tasks via Python programming.

On a variety of English evaluation tasks, DeepSeekMath-Base demonstrates performance on par with the proprietary Minerva 540B \citep{minerva}, and consistently outperforms all open-source base models—such as Mistral 7B \citep{mistral} and Llemma-34B \citep{llemma}—whether or not these alternatives have received math-specific pre-training, often by a wide margin. Remarkably, DeepSeekMath-Base achieves even higher results on Chinese assessment benchmarks, which may be attributed to our inclusion of high-quality, non-English math pre-training data, as opposed to the English-only datasets used in prior research \citep{minerva,llemma}. Through the application of mathematical instruction fine-tuning and reinforcement learning, the extended models DeepSeekMath-Instruct and DeepSeekMath-RL attain outstanding results, exceeding 50\% accuracy on the highly challenging MATH dataset—marking a first for open-source systems.

\item \textbf{Formal Mathematics}: The capability of DeepSeekMath-Base is assessed through the informal-to-formal theorem proving benchmark presented in \citep{dsp_proof}, utilizing the miniF2F dataset \citep{minif2f} with the Isabelle proof assistant \citep{isabelle}. Results show that DeepSeekMath-Base achieves impressive performance in few-shot settings for autoformalization.

\item \textbf{Natural Language Understanding, Reasoning, and Code}: To obtain a holistic evaluation of the model’s proficiency in general comprehension, logical reasoning, and programming, we test DeepSeekMath-Base on the Massive Multitask Language Understanding (MMLU) benchmark \citep{mmlu}, which spans 57 multiple-choice tests across a variety of fields, as well as on BIG-Bench Hard (BBH) \citep{bbh}, a suite of 23 intricate tasks primarily focused on multi-step reasoning. For assessing programming ability, we use HumanEval \citep{codex} and MBPP \citep{mbpp}, both established standards for code model evaluation. Pre-training on mathematics has been found to enhance the model's performance in language understanding and complex reasoning tasks.

\section{Math Pre-Training}

\subsection{Data Collection and Decontamination} This section describes the methodology employed to build the DeepSeekMath Corpus utilizing data sourced from Common Crawl. Illustrated in Figure \ref{fig:data_collect}, we introduce an iterative workflow designed to efficiently assemble a substantial mathematical dataset from Common Crawl, beginning with an initial seed corpus comprised of a limited yet highly curated math-focused dataset. It should be emphasized that this strategy is transferable and can be effectively used for corpus construction in other fields, including coding.

We begin our process by selecting OpenWebMath \citep{openwebmath}—a curated dataset of superior mathematical content sourced from the web—as our seed corpus. This corpus serves as the foundation for training a fastText model \citep{joulin2016fasttext} designed to identify additional web pages with mathematical content akin to OpenWebMath. For training, we sample 500,000 entries from the seed dataset to act as positive instances, while an additional 500,000 web pages randomly drawn from Common Crawl provide the negative class. Training is conducted using an open-source implementation\footnote{\url{https://fasttext.cc}}, setting hyperparameters as follows: a vector size of 256, learning rate of 0.1, a maximum word n-gram length of 3, a minimum word frequency threshold of 3, and 3 training epochs. To curtail the scale of the Common Crawl dataset, we utilize URL deduplication and near-duplicate removal, ultimately generating a corpus of 40B distinct HTML pages. Applying the trained fastText model to this deduplicated dataset, we extract pages likely to contain mathematical content. We further mitigate low-quality selections by ranking these candidates based on fastText’s predicted scores, preserving only those with the highest relevance. We assess the preserved data volume through pre-training trials at thresholds of the top 40B, 80B, 120B, and 160B tokens. For the initial round, we opt to retain pages comprising the top 40B tokens.

Following the initial round of data acquisition, a significant number of mathematical web pages are not captured, primarily due to the fastText model's reliance on a set of positive examples that are not sufficiently varied. To address this, we expand the seed corpus by incorporating further sources of mathematical web content, aiming to refine the fastText model’s efficacy. The procedure begins with segmenting the entire Common Crawl dataset by domain, where a domain refers to web pages originating from an identical base URL. For each domain, we determine the proportion of its web pages that were captured in the first data collection iteration. Domains with more than 10\% of their pages included in the initial sweep are designated as math-related (such as \url{mathoverflow.net}). Next, we conduct manual annotation of URLs within these domains that are associated with mathematical material (for instance, \url{mathoverflow.net/questions}). Any remaining web pages under these annotated URLs that were not collected previously are added to the growing seed corpus. This expanded set of positive examples enhances the training of the fastText model, which, in turn, increases recall of mathematical content in further iterations. Across four data collection cycles, this process yields 35.5M mathematical web pages, encompassing 120B tokens in total. By the fourth iteration, it becomes apparent that approximately 98\% of the mathematical data was already acquired by the third iteration, prompting the decision to discontinue further data collection.

In order to prevent contamination of our benchmarks, we adopt the filtering procedure outlined by \cite{deepseek-coder}: web pages that include questions or answers originating from major English mathematical benchmarks such as GSM8K \citep{gsm8k} and MATH \citep{MATH}, as well as Chinese benchmarks like CMATH \citep{wei2023cmath} and AGIEval \citep{agieval}, are removed. Specifically, we eliminate from our mathematical training data any text segment that contains a 10-gram phrase exactly matching a substring from any of these benchmark datasets. For benchmark contents whose length falls between 3 and 9 grams, we apply an exact match strategy to further filter potentially contaminated sources from the web pages.

\subsection{Validating the Quality of the DeepSeekMath~Corpus}

We run pre-training experiments to investigate how the DeepSeekMath~Corpus is compared with the recently released math-training corpora:

\begin{itemize}[topsep=0pt]
    \item \textbf{MathPile} \citep{mathpile}: an 8.9B token corpus derived from multiple sources including textbooks, Wikipedia, ProofWiki, CommonCrawl, StackExchange, and arXiv, with arXiv comprising more than 85\% of the material;
    \item \textbf{OpenWebMath} \citep{openwebmath}: a 13.6B token dataset generated by extracting mathematics-related data from CommonCrawl;
    \item \textbf{Proof-Pile-2} \citep{llemma}: a collection tailored to mathematics that combines OpenWebMath, AlgebraicStack (accounting for 10.3B tokens in mathematical code), and arXiv manuscripts (amounting to 28.0B tokens). Our usage of Proof-Pile-2 adheres to the arXiv:Web:Code proportion of 2:4:1, as proposed in \cite{llemma}.
\end{itemize}

\subsubsection{Training Setting}
\label{sec:quality-policy}
We perform mathematical pre-training on a general-purpose language model comprising 1.3 billion parameters, architecturally consistent with DeepSeek LLMs \citep{deepseek-llm}, referred to here as DeepSeek-LLM 1.3B. The model is fine-tuned independently on each math dataset for a total of 150 billion tokens. Training is executed utilizing the resource-efficient HAI-LLM \citep{haillm} framework. In alignment with the methodology established for DeepSeek LLMs, we employ the AdamW optimizer \citep{adamW}, setting $\beta_1=0.9$, $\beta_2=0.95$, and $\mathrm{weight\_decay}=0.1$. A multi-stage learning rate scheduling strategy is adopted: the learning rate is warmed up to its peak over the initial 2{,}000 steps, then reduced to 31.6\% of its maximum by the 80\% training mark, and further annealed to 10.0\% of the peak by 90\% completion. The highest learning rate is capped at 5.3e-4. Training is carried out using a batch size of 4 million tokens, with a maximum sequence length (context length) of 4,000 tokens.

\subsubsection{Evaluation Results}

\textbf{The DeepSeekMath~Corpus is of high quality, covers multilingual mathematical content, and is the largest in size.}

\begin{itemize}[topsep=0pt]
    \item \textbf{High-quality}:
    We assess downstream task performance across 8 mathematical evaluation sets utilizing few-shot chain-of-thought prompting \cite{cot}. Table \ref{tab:corpora_comparison} highlights the superior outcomes achieved by the model trained with the DeepSeekMath~Corpus. In Figure \ref{fig:corpora_comparisons}, models developed on DeepSeekMath Corpus surpass those trained with Proof-Pile-2 after exposure to 50B tokens (equivalent to one full epoch of Proof-Pile-2), reflecting the consistently higher average data quality of DeepSeekMath Corpus.
    \item \textbf{Multilingual}:
    The DeepSeekMath~Corpus contains material spanning various languages, with English and Chinese as its principal components. Table \ref{tab:corpora_comparison} demonstrates that employing DeepSeekMath~Corpus in training notably improves mathematical reasoning in both English and Chinese. In comparison, other available mathematical datasets, which are predominantly focused on English, exhibit limited benefits or can even adversely affect performance for Chinese mathematical reasoning.
    \item \textbf{Large-scale}:
    The scale of the DeepSeekMath~Corpus substantially exceeds that of other mathematical corpora. Figure \ref{fig:corpora_comparisons} illustrates that DeepSeek-LLM 1.3B, once trained on DeepSeekMath~Corpus, achieves a much steeper improvement trajectory and experiences more durable gains. By contrast, baseline corpora are considerably smaller in size, have undergone multiple repeat passes during training, and their model performance plateaus rapidly as a result.
\end{itemize}

\subsection{Training and Evaluating DeepSeekMath-Base 7B}

This section presents DeepSeekMath-Base 7B, a foundational model characterized by robust mathematical reasoning capabilities. The model builds upon DeepSeek-Coder-Base-v1.5 7B \citep{deepseek-coder} and undergoes training on 500B tokens. The training data composition is as follows: 56\% originates from the DeepSeekMath Corpus, 4\% from AlgebraicStack, 10\% from arXiv, 20\% comprises Github code, while the final 10\% consists of natural language samples from Common Crawl in both English and Chinese. With regard to the training setup, we largely adhere to the approach detailed in Section \ref{sec:quality-policy}, with two key adjustments: the peak learning rate is set to 4.2e-4, and the batch size is fixed at 10M tokens.

We undertake an extensive evaluation of DeepSeekMath-Base 7B’s proficiency in mathematics, examining its capacity to generate complete mathematical solutions independently, tackle mathematical tasks utilizing auxiliary tools, and execute formal theorem-proving procedures. In addition to this mathematical focus, we further present a broader characterization of the base model by assessing its abilities in natural language comprehension, logical reasoning, and coding.

\paragraph{Mathematical Problem Solving with Step-by-Step Reasoning}

We assess the ability of DeepSeekMath-Base to solve mathematical questions using few-shot chain-of-thought prompting \citep{cot}, evaluating across eight benchmarks in both English and Chinese. These benchmarks include those focused on quantitative reasoning—such as GSM8K \citep{gsm8k}, MATH \citep{MATH}, and CMATH \citep{wei2023cmath}—as well as multiple-choice tasks, like MMLU-STEM \citep{mmlu} and Gaokao-MathQA \citep{agieval}. Together, they span a wide spectrum of mathematical domains, from elementary arithmetic to advanced college-level topics.

As indicated in Table \ref{tab:base_math_cot}, DeepSeekMath-Base 7B demonstrates superior results across all eight evaluated benchmarks when compared to other open-source base models, such as the prominent general-purpose Mistral 7B \citep{mistral} and the more recently introduced Llemma 34B \citep{llemma}, which received specialized mathematical training on Proof-Pile-2 \citep{llemma}. In particular, on the highly challenging MATH dataset, DeepSeekMath-Base achieves an absolute lead of more than 10\% over other open-source base models. Furthermore, it also outperforms Minerva 540B \citep{minerva}—a substantially larger closed-source model, based on PaLM \citep{palm} and enhanced through further mathematical corpus training—despite Minerva’s size advantage of approximately 77 times more parameters.

\paragraph{Mathematical Problem Solving with Tool Use} Our assessment of program-assisted mathematical reasoning is carried out on the GSM8K and MATH datasets employing few-shot program-of-thought prompting \citep{pot,pal}. The models receive prompts to address each problem by generating a Python script, making use of libraries like \textit{math} and \textit{sympy} to perform complex computations where necessary. The output produced upon running the generated code is treated as the solution. As illustrated in Table \ref{tab:base_math_pot_proof}, DeepSeekMath-Base 7B achieves better results than the previously leading Llemma 34B model.

\paragraph{Formal Mathematics}

Automating formal proofs plays a crucial role in maintaining the correctness and dependability of mathematical arguments while improving productivity—a topic that has received growing interest in recent years. We assess the capabilities of DeepSeekMath-Base 7B on the informal-to-formal proving task from \citep{dsp_proof}, where the objective is to construct a formal proof given an informal proposition, its formal analogue, and an informal reasoning. For our experiments, we utilize miniF2F \citep{minif2f}, a standardized benchmark for formalizing Olympiad-level mathematical challenges, and prompt the generation of Isabelle formal proofs using few-shot examples for each instance. In line with \cite{dsp_proof}, we use the model to create high-level proof outlines and subsequently employ the automated prover Sledgehammer \citep{sledgehammer} to elaborate the remaining proof steps. As depicted in Table \ref{tab:base_math_pot_proof}, DeepSeekMath-Base 7B achieves notable results in the area of proof autoformalization.

\paragraph{Natural Language Understanding, Reasoning, and Code}

We assess natural language understanding using MMLU \citep{mmlu}, evaluate reasoning ability through BBH \citep{bbh}, and test programming proficiency on HumanEval \citep{codex} and MBPP \citep{mbpp}. Table \ref{tab:understanding_reasoning_code} demonstrates that DeepSeekMath-Base 7B achieves marked improvements over its predecessor, DeepSeek-Coder-Base-v1.5 \citep{deepseek-coder}, in both MMLU and BBH, highlighting the beneficial effect of mathematical training on comprehension and reasoning tasks. Furthermore, with continued exposure to code tokens, DeepSeekMath-Base 7B preserves the coding benchmark performance observed in DeepSeek-Coder-Base-v1.5. Collectively, DeepSeekMath-Base 7B delivers substantially better results than the general-purpose Mistral 7B \citep{mistral} on all three reasoning and programming evaluations.

\section{Supervised Fine-Tuning}

\subsection{SFT Data Curation}

We develop an instruction-tuning dataset for mathematics that encompasses both English and Chinese questions, sampling from a variety of mathematical disciplines and complexity tiers. Each problem is matched with a corresponding solution formatted as chain-of-thought (CoT) \citep{cot}, program-of-thought (PoT) \citep{pot,pal}, or through tool-integrated reasoning \citep{tora}. The training corpus contains a total of 776K examples.

\begin{itemize}[topsep=0pt]
    \item \textbf{English mathematical datasets}: We provide tool-integrated solution annotations for problems from GSM8K and MATH, in addition to leveraging a sample of MathInstruct \citep{MathInstruct} and the Lila-OOD training data \citep{lila}, both of which contain problems addressed via CoT or PoT strategies. The English dataset encompasses a wide range of mathematical disciplines, including algebra, calculus, number theory, probability, and geometry.
    \item \textbf{Chinese mathematical datasets}: Our compilation consists of Chinese K-12 mathematics questions that cover 76 distinct sub-topics, such as linear equations, with each problem featuring annotated solutions in both CoT and tool-integrated reasoning styles.
\end{itemize}

\subsection{Training and Evaluating DeepSeekMath-Instruct 7B}

This section presents DeepSeekMath-Instruct 7B, a model refined through mathematical instruction tuning originating from DeepSeekMath-Base. To prepare training data, examples are joined together in random order, with the total not exceeding a 4K token context window. The training procedure runs for 500 iterations, employing a batch size of 256 and maintaining a fixed learning rate of 5e-5 throughout.

We evaluate models' mathematical performance both without and with tool use, on 4 quantitative reasoning benchmarks in English and Chinese. We benchmark our model against the leading models of the time:

\begin{itemize}[topsep=0pt]
    \item \textbf{Closed-source models} comprise the following: (1) the GPT series, where GPT-4 \citep{gpt4} and GPT-4 Code Interpreter~\footnote{\url{https://openai.com/blog/chatgpt-plugins\#\#code-interpreter}} represent the most advanced variants, (2) Gemini Ultra and Pro \citep{gemini}, (3) Inflection-2 \citep{inflection-2}, (4) Grok-1~\footnote{\url{https://x.ai/model-card}}, as well as several models made available by Chinese tech firms, such as (5) Baichuan-3~\footnote{\url{https://www.baichuan-ai.com}}, and (6) the latest version of GLM, GLM-4~\footnote{\url{https://open.bigmodel.cn/dev/api\#glm-4}}.

    from the GLM family \citep{glm}.

These models are designed for broad, general-purpose tasks and have typically been subjected to various alignment protocols. 

\item \textbf{Open-source models} feature a range of architectures: standard general-purpose models such as (1) DeepSeek-LLM-Chat 67B \citep{deepseek-llm}, (2) Qwen 72B \citep{qwen}, (3) SeaLLM-v2 7B \citep{seallm}, and (4) ChatGLM3 6B \citep{chatglm3}. In addition, models optimized for mathematics include (5) InternLM2-Math 20B~\footnote{\url{https://github.com/InternLM/InternLM-Math}}, an extension of InternLM2 with additional training on mathematical datasets and further refined with instruction tuning; (6) Math-Shepherd-Mistral 7B, which involves applying PPO training \citep{schulman2017proximal} to Mistral 7B \citep{mistral} using a reward model that leverages process supervision; (7) the WizardMath suite \citep{wizardmath}, which enhances the mathematical capabilities of Mistral 7B and Llama-2 70B \citep{llama2} via evolve-instruct (an instruction tuning technique using AI-generated instructions) and PPO, mainly training on problems from GSM8K and MATH; (8) MetaMath 70B \citep{metamath}, which consists of Llama-2 70B fine-tuned on enriched GSM8K and MATH datasets; (9) ToRA 34B \cite{tora}, a version of CodeLlama 34B adapted through fine-tuning for mathematical reasoning that incorporates tool use; and (10) MAmmoTH 70B \citep{MathInstruct}, which applies instruction tuning to Llama-2 70B using the MathInstruct data.

As indicated in Table \ref{tab:sft_rl_math}, when tool use is not permitted during evaluation, DeepSeekMath-Instruct 7B delivers robust step-by-step reasoning capabilities. Remarkably, on the challenging competition-level MATH dataset, this model outperforms all other open-source alternatives and most commercial solutions (including Inflection-2 and Gemini Pro) by a margin of at least 9\% in absolute terms. This performance advantage holds even compared to models with far greater parameter counts (such as Qwen 72B) or those specifically optimized via math-oriented reinforcement learning techniques (such as WizardMath-v1.1 7B). Although DeepSeekMath-Instruct is competitive with Chinese proprietary models like GLM-4 and Baichuan-3 on MATH, it does not yet reach the performance levels of GPT-4 or Gemini Ultra.

When assessed in a context permitting both natural language reasoning and programmatic tool application for problem-solving, DeepSeekMath-Instruct 7B attains nearly 60\% accuracy on MATH, outperforming all previously released open-source models. Across additional benchmarks, our model’s performance rivals that of DeepSeek-LLM-Chat 67B, the former leading model in this area, despite its much larger parameter size (10 times greater).

\subsection{Group Relative Policy Optimization} Reinforcement learning (RL) has shown considerable promise in enhancing the mathematical reasoning capabilities of LLMs following the Supervised Fine-Tuning (SFT) phase \citep{wang2023math,wizardmath}. Here, we present Group Relative Policy Optimization (GRPO), our novel RL algorithm designed for both efficiency and effectiveness.

\subsubsection{From PPO to GRPO} Proximal Policy Optimization (PPO) \citep{schulman2017proximal} is an actor-critic RL algorithm that is widely used in the RL fine-tuning stage of LLMs \citep{ouyang2022training}. In particular, it optimizes LLMs by maximizing the following surrogate objective:

\begin{equation}
\footnotesize
    \mathcal{J}_{PPO}(\theta) = \mathbb{E}{[q \sim P(Q), o \sim \pi_{\theta_{old}}(O|q)]} \frac{1}{|o|} \sum_{t=1}^{|o|} \min \left[ \frac{\pi_\theta(o_{t} | q, o_{<t})}{\pi_{\theta_{old}}(o_{t} | q, o_{<t})} A_{t}, \text{clip} \left( \frac{\pi_\theta(o_{t} | q, o_{<t})}{\pi_{\theta_{old}}(o_{t} | q, o_{<t})}, 1 - \varepsilon, 1 + \varepsilon \right)  A_{t} \right] ,
\end{equation}

The objective function for PPO, denoted as $\mathcal{J}_{PPO}(\theta)$, involves taking the expectation with respect to $q$ drawn from $P(Q)$ and $o$ sampled according to the prior policy $\pi_{\theta_{old}}(O|q)$. For each token position $t$ in the sequence $o$, the term evaluates the minimum between the product of the advantage $A_t$ and the policy likelihood ratio $\frac{\pi_\theta(o_t | q, o_{<t})}{\pi_{\theta_{old}}(o_t | q, o_{<t})}$, and the same product after the ratio is clipped within $[1-\varepsilon, 1+\varepsilon]$ to enforce a trust region. The sum is averaged across sequence length $|o|$.

Here, $\pi_{\theta}$ denotes the current policy while $\pi_{\theta_{old}}$ represents the previous policy model. The variables $q$ and $o$ refer to the sampled questions and their associated outputs, obtained from the question corpus and generated using the old policy $\pi_{\theta_{old}}$. The hyper-parameter $\varepsilon$ serves as a clipping threshold in PPO, which assists in enhancing the stability of the optimization process. $A_t$ designates the advantage, calculated via Generalized Advantage Estimation (GAE) \citep{gae}, utilizing the observed rewards $\{r_{\ge t}\}$ in conjunction with a learned value function $V_{\psi}$. Consequently, PPO involves concurrently training a value function and the policy network. To address excessive reward optimization, it is customary to include a per-token KL-divergence penalty—derived from a reference model—within the reward computation at each token, following the methodology of \citep{ouyang2022training}, i.e.,

\begin{equation}
    r_{t} = r_\varphi(q, o_{\le t}) - \beta \log\frac{\pi_{\theta}(o_{t}|q, o_{<t})}{\pi_{ref}(o_{t}|q, o_{<t})},
\label{eq:PPO-reward}
\end{equation}

Here, $r_\varphi$ denotes the reward model, $\pi_{ref}$ represents the reference model—commonly set as the initial SFT model—and $\beta$ serves as the scaling factor for the KL divergence penalty.

As the value function employed in PPO is typically another model of comparable size as the policy model, it brings a substantial memory and computational burden. Additionally, during RL training, the value function is treated as a baseline in the calculation of the advantage for variance reduction. While in the LLM context, usually only the last token is assigned a reward score by the reward model, which may complicate the training of a value function that is accurate at each token. To address this, as shown in Figure \ref{fig:grpo}, we propose Group Relative Policy Optimization (GRPO), which obviates the need for additional value function approximation as in PPO, and instead uses the average reward of multiple sampled outputs, produced in response to the same question, as the baseline. More specifically, for each question $q$, GRPO samples a group of outputs $\{o_1, o_2, \cdots, o_G\}$  from the old policy  $\pi_{\theta_{old}}$  and then optimizes the policy model by maximizing the following objective:

\begin{equation}
\footnotesize
\begin{split}
    \mathcal{J}_{GRPO}(\theta) &= \mathbb{E}{[q \sim P(Q), \{o_i\}_{i=1}^G \sim \pi_{\theta_{old}}(O|q)]}  \\
    & \frac{1}{G} \sum_{i=1}^G \frac{1}{|o_i|} \sum_{t=1}^{|o_i|} \left\{ \min \left[ \frac{\pi_\theta(o_{i,t} | q, o_{i,<t})}{\pi_{\theta_{old}}(o_{i,t} | q, o_{i,<t})} \hat{A}_{i,t}, \ \text{clip}\left( \frac{\pi_\theta(o_{i,t} | q, o_{i,<t})}{\pi_{\theta_{old}}(o_{i,t} | q, o_{i,<t})}, 1 - \varepsilon, 1 + \varepsilon \right) \hat{A}_{i,t} \right] - \beta \mathbb{D}_{KL}\left[\pi_{\theta} || \pi_{ref}\right] \right\} ,
\end{split}
\label{eq:GRPO-obj}
\end{equation}

where $\varepsilon$ and $\beta$ are hyper-parameters, and $\hat{A}_{i,t}$  is the advantage calculated based on relative rewards of the outputs inside each group only, which will be detailed in the following subsections. The group relative way that GRPO leverages to calculate the advantages, aligns well with the comparative nature of rewards models,  as reward models are typically trained on datasets of comparisons between outputs on the same question. Also note that, instead of adding KL penalty in the reward, GRPO regularizes by directly adding the KL divergence between the trained policy and the reference policy to the loss, avoiding complicating the calculation of  $\hat{A}_{i,t}$. And different from the KL penalty term used in (\ref{eq:PPO-reward}), we estimate the KL divergence with the following unbiased estimator \citep{kl_approx}:

\begin{equation}
\small
    \mathbb{D}_{KL}\left[\pi_{\theta} || \pi_{ref}\right] = \frac{\pi_{ref}(o_{i,t}|q,o_{i,<t})}{\pi_{\theta}(o_{i,t}|q,o_{i,<t})}- \log\frac{\pi_{ref}(o_{i,t}|q,o_{i,<t})}{\pi_{\theta}(o_{i,t}|q,o_{i,<t})} - 1,
\end{equation}

which is guaranteed to be positive.

\begin{algorithm}[t]
  \small
  \caption{Iterative Group Relative Policy Optimization}
  \textbf{Input} initial policy model $\pi_{\theta_{\text{init}}}$; reward models $r_\varphi$; task prompts $\mathcal{D}$; 
  hyperparameters $\varepsilon$, $\beta$, $\mu$
  \begin{algorithmic}[1]
    \State Initialize the policy $\pi_\theta$ as $\pi_{\theta_{\text{init}}}$
    \For{iteration = 1, \dots, I}
       \State Set reference policy $\pi_{ref}$ to the current model: $\pi_{ref} \leftarrow \pi_{\theta}$
      \For{step = 1, \dots, M}
      \State Randomly draw a minibatch $\mathcal{D}_b$ from $\mathcal{D}$
      \State Save current parameters as previous: $\pi_{\theta_{old}} \leftarrow \pi_{\theta}$ 
      \State For every prompt $q$ in $\mathcal{D}_b$, generate $G$ samples $\{o_i\}_{i=1}^G \sim \pi_{\theta_{old}} (\cdot \mid q)$
      \State For each candidate response $o_i$, compute reward values $\{r_i\}_{i=1}^{G}$ using $r_\varphi$
      \State Perform group relative advantage estimation to obtain $\hat{A}_{i,t}$ for each $t$-th token in $o_i$
      \For{GRPO iteration = 1, \dots, $\mu$}
        \State Adjust policy parameters $\pi_{\theta}$ to maximize the GRPO target in Equation \ref{eq:GC-GRPO}
      \EndFor
    \EndFor
    \State Continuously refine $r_\varphi$ utilizing replay-based training. 
    \EndFor 
  \end{algorithmic}
  \textbf{Output} $\pi_\theta$
  \label{alg:iter-grpo}
\end{algorithm}

\subsubsection{Outcome Supervision RL with GRPO} Specifically, for each question $q$, a set of candidate outputs $\{o_1, o_2, \cdots, o_G\}$ is generated via sampling from the prior policy $\pi_{\theta_{old}}$. Each output is then evaluated with a reward model, resulting in $G$ corresponding reward values $\mathbf{r} = \{r_1, r_2, \cdots, r_G\}$. To normalize these reward signals, the group mean is subtracted from each reward, and the result is divided by the standard deviation of the group. Outcome supervision assigns this normalized reward at the conclusion of each output $o_i$, attributing it as the advantage $\hat{A}_{i,t}$ for every token within the output, formally: $\hat{A}_{i, t} = \widetilde{r}_i = \frac{r_i- {\rm mean}(\mathbf{r})}{{\rm std}(\mathbf{r})}$. Finally, policy optimization proceeds by maximizing the objective specified in equation (\ref{eq:GRPO-obj}).

\subsubsection{Process Supervision RL with GRPO} Outcome supervision restricts feedback to a single reward delivered at the conclusion of each output, which can be inadequate for efficiently guiding policies in complex mathematical reasoning tasks. Building on the approach from \cite{wang2023math}, we investigate process supervision, where rewards are assigned at each intermediate reasoning step. Concretely, for a question $q$ and a set of $G$ sampled outputs $\{o_1, o_2, \cdots, o_G\}$, a process reward model evaluates every step in the outputs and assigns rewards accordingly: $\mathbf{R} = \{ \{r_1^{index(1)},\cdots,r_1^{index(K_1)}\}, \cdots,  \{r_G^{index(1)},\cdots,r_G^{index(K_G)}\} \}$, where $index(j)$ indicates the token index at which the $j$-th reasoning step ends, and $K_i$ denotes the number of steps in the $i$-th output. To ensure comparability across reward magnitudes, the obtained stepwise rewards are standardized using their mean and standard deviation, following $\widetilde{r}_i^{index(j)} = \frac{r_i^{index(j)} - {\rm mean(\mathbf{R})}}{{\rm std(\mathbf{R})}}$. Afterwards, the process supervision computes the token-wise advantage as the accumulated sum of all subsequent normalized rewards from that token onward, specifically $\hat{A}_{i, t} = \sum_{index(j) \ge t} \widetilde{r}_i^{index(j)}$. The policy is then updated to maximize the objective set out in equation (\ref{eq:GRPO-obj}).

\subsubsection{Iterative RL with GRPO} Over the course of reinforcement learning, the existing reward model may become inadequate for guiding the evolving policy model. To address this limitation, we investigate an iterative RL framework utilizing GRPO. As depicted in Algorithm \ref{alg:iter-grpo}, the iterative GRPO approach involves creating updated training datasets for the reward model using newly sampled outputs from the policy model. The reward model is incrementally updated via a replay strategy that retains 10\% of previous training samples. Subsequently, the reference model is synchronized with the latest policy model, and further updates of the policy model proceed using the enhanced reward model.

\subsection{Training and Evaluating DeepSeekMath-RL}

Our reinforcement learning experiments are centered on DeepSeekMath-Instruct 7B. For RL training, we use chain-of-thought questions derived from GSM8K and MATH within the SFT dataset, amounting to approximately 144K samples. To specifically analyze the effects of RL when target benchmarks are unseen during training, we omit SFT questions unrelated to these tasks. The reward model’s training set is assembled in accordance with the approach outlined by \citep{wang2023math}. We initialize the reward model using DeepSeekMath-Base 7B, employing a learning rate of 2e-5. For GRPO, the policy model operates with a learning rate of 1e-6, and the KL penalty is set to 0.04. For every question, $64$ completions are generated. The sequence length limit is set at 1024 tokens, and we use a batch size of 1024 during training. After each exploration iteration, the policy model undergoes a single parameter update. To assess model performance, DeepSeekMath-RL 7B is evaluated using the same benchmarks as DeepSeekMath-Instruct 7B. Within this framework, GSM8K and MATH tasks requiring chain-of-thought reasoning are considered in-domain, while all other benchmarks represent out-of-domain challenges.

Table \ref{tab:sft_rl_math} presents a comparison of the results for both open-source and proprietary models, evaluated using chain-of-thought and tool-based reasoning on benchmarks in English and Chinese. The findings reveal that: 1) DeepSeekMath-RL 7B achieves 88.2\% accuracy on GSM8K and 51.7\% on MATH when employing chain-of-thought reasoning. These outcomes outperform every other open-source model in the 7B to 70B parameter range and also exceed those of most closed-source alternatives. 2) Notably, the training of DeepSeekMath-RL 7B is confined solely to instruction tuning data in chain-of-thought format from GSM8K and MATH, with initialization from DeepSeekMath-Instruct 7B. Despite this limitation in training data, DeepSeekMath-RL 7B consistently surpasses DeepSeekMath-Instruct 7B in all evaluation criteria, providing evidence for the benefit of incorporating reinforcement learning.

\section{Discussion}
Here, we present the results observed in our investigations of both pre-training and reinforcement learning (RL) experiments.

\subsection{Lessons Learnt in Pre-Training}

We begin by detailing our observations during the pre-training phase. Throughout this section, unless stated otherwise, we follow the training configuration described in Section \ref{sec:quality-policy}. Additionally, for references to the DeepSeekMath Corpus herein, we employ an 89B-token dataset corresponding to the second round of data gathering.

\subsubsection{Code Training Benefits Mathematical Reasoning}

An often-cited but largely untested hypothesis posits that exposure to code can enhance reasoning skills. In this work, we seek to partially address this claim in the context of mathematics, demonstrating that code training enhances models’ proficiency in mathematical reasoning tasks, regardless of whether external tools are employed.

To study how code training affects mathematical reasoning, we experimented with the following two-stage training and one-stage training settings:

\textbf{Two-Stage Training}

\begin{itemize}[topsep=0pt]
    \item \textbf{Code Training for 400B Tokens $\rightarrow$ Math Training for 150B Tokens}: The DeepSeek-LLM 1.3B model is initially pretrained on 400 billion tokens sourced from code, and subsequently fine-tuned with 150 billion tokens of mathematical content;
    \item \textbf{General Training for 400B Tokens $\rightarrow$ Math Training for 150B Tokens}: For comparison, a variant is trained on 400 billion tokens selected from a comprehensive general-purpose corpus developed by DeepSeek-AI during the first phase, followed by the same math-focused fine-tuning, aiming to evaluate whether pretraining on code tokens confers a benefit over general tokens for mathematical reasoning tasks.
\end{itemize}

\textbf{One-Stage Training}

\begin{itemize}[topsep=0pt]
    \item \textbf{Math Training Utilizing 150B Tokens}: DeepSeek-LLM 1.3B undergoes training on a dataset comprising 150B mathematical tokens;
    \item \textbf{Joint Training on 400B Code Tokens and 150B Math Tokens}: When mathematical instruction is introduced after code training, there is a decline in code task performance. We examine if simultaneously training on both code and math tokens in a unified stage can enhance mathematical reasoning capabilities and help prevent catastrophic forgetting.
\end{itemize}

\paragraph{Results}
Tables \ref{tab:code-math} and \ref{tab:code-math-general-eval} present the downstream results obtained across various training configurations.

Incorporating code training facilitates improvements in program-aided mathematical reasoning tasks, regardless of whether a two-stage or one-stage training approach is adopted. As indicated in Table \ref{tab:code-math}, with the two-stage training paradigm, code training on its own brings about substantial gains in the model’s capability to address GSM8K and MATH tasks using Python. Supplementary math training in a subsequent stage leads to additional advancements. Notably, when utilizing a one-stage training framework, where both code and math tokens are interleaved, the approach not only alleviates the problem of catastrophic forgetting seen in two-stage training, but also enhances both general coding abilities (see Table \ref{tab:code-math-general-eval}) and competence in program-supported mathematical reasoning (see Table \ref{tab:code-math}).

Training on code additionally enhances mathematical reasoning skills even in the absence of external tools. Within the two-stage training framework, the first phase of code-focused instruction yields noticeable, though not maximal, improvements. Furthermore, this phase facilitates greater efficiency during subsequent mathematical training, cumulatively producing superior outcomes. In contrast, employing a single-stage training approach that simultaneously introduces both code and mathematical tokens appears to diminish reasoning capabilities when tools are not utilized. This may be attributable to the restricted size of DeepSeek-LLM 1.3B, which likely constrains its ability to adequately integrate both coding and mathematical information in a single training period.

\subsubsection{ArXiv Papers Seem Ineffective in Improving Mathematical Reasoning}

ArXiv papers are commonly included as a component of math pre-training data \citep{minerva,gpt-f,llemma,mathpile}. However, detailed analysis regarding their impact on mathematical reasoning has not been extensively conducted. Perhaps counter-intuitively, according to our experiments, arXiv papers seem ineffective in improving mathematical reasoning. We experiment with models of different sizes, including DeepSeek-LLM 1.3B and DeepSeek-Coder-Base-v1.5 7B \citep{deepseek-coder}, using arXiv corpora that underwent varied processing pipelines:

\begin{itemize}[topsep=0pt]
    \item \textbf{MathPile} \citep{mathpile}: a dataset containing 8.9B tokens, primarily composed of scientific arXiv articles (over 85\%) and curated through specific cleaning and filtering heuristics;
    \item \textbf{ArXiv-RedPajama} \citep{redpajama}: comprises all arXiv LaTeX source files with extraneous elements such as preambles, comments, macros, and bibliographic entries excluded, resulting in a 28.0B token corpus.
\end{itemize}

For our experimental setup, we independently pretrain DeepSeek-LLM 1.3B with 150B tokens and DeepSeek-Coder-Base-v1.5 7B with 40B tokens on each distinct arXiv dataset. The findings indicate that leveraging arXiv articles alone does not meaningfully enhance mathematical reasoning capabilities. When models are exposed solely to arXiv-based corpora, neither exhibits substantial progress—and in some cases, even regress—across a range of mathematical evaluation suites with varying degrees of difficulty considered in this work. The evaluation encompasses datasets focused on quantitative reasoning such as GSM8K and MATH (see Table \ref{tab:arxiv-cot}), multiple-choice assessments like MMLU-STEM (Table \ref{tab:arxiv-cot}), as well as formal mathematics benchmarks, including miniF2F (Table \ref{tab:arxiv_atp}).

However, this conclusion has its limitations and should be taken with a grain of salt. We have not yet studied:

\begin{itemize}[topsep=0pt]
    \item How arXiv tokens influence particular mathematical activities beyond the scope of this work, for instance, transforming formal theorems or proofs into their informal counterparts;
    \item The interaction between arXiv tokens and additional data modalities or sources;
    \item If scaling models up would enhance or reveal the advantages associated with utilizing arXiv papers.
\end{itemize}

Thus, further exploration is required, which we leave for future studies.

\subsection{Insights of Reinforcement Learning}
\subsubsection{Towards to a Unified Paradigm} Here, we introduce a comprehensive framework for examining a range of training techniques, including SFT, RFT, DPO, PPO, and GRPO. We also perform empirical studies to investigate the elements within this unified framework. In general terms, the gradient of the objective function with respect to the parameter $\theta$ for a particular training approach may be expressed as:

\begin{equation}
    \nabla_{\theta}\mathcal{J}_{\textcolor{red}{\mathcal{A}}}(\theta) = \mathbb{E}[\underbrace{(q,o) \sim \textcolor{red}{\mathcal{D}}}_{Data \ Source}]\left( \frac{1}{|o|} \sum_{t=1}^{|o|}  \underbrace{GC_{{\mathcal{A}}}(q, o, t, \textcolor{red}{\pi_{{rf}}})}_{Gradient \ Coefficient}  \nabla_{\theta}\log \pi_{\theta}(o_t | q, o_{<t})\right).
\label{eq:objective}
\end{equation}

There exist three key components: 1) \textit{Data Source $\mathcal{D}$}, which determines the training data; 2) \textit{Reward Function $\pi_{{rf}}$}, which is the source of the training reward signal; 3) \textit{Algorithm $\mathcal{A}$}: which processes the training data and the reward signal to the gradient coefficient $GC$ that determines the magnitude of the penalty or reinforcement for the data. We analyze several representative methods based on such a unified paradigm:

\begin{itemize}
    \item \textbf{Supervised Fine-tuning (SFT)}: This approach involves further training a pretrained model on a dataset curated by humans specifically for SFT, enabling the model to align more closely with human preferences.
    \item \textbf{Rejection Sampling Fine-tuning (RFT)}: In RFT, the SFT-trained model undergoes additional fine-tuning on a subset of its own generated outputs. These outputs, produced in response to SFT prompts, are filtered such that only those with accurate answers are retained for training.
    \item \textbf{Direct Preference Optimization (DPO)}: DPO enhances the SFT-trained model by applying fine-tuning on an expanded set of outputs that the model produces, utilizing a pairwise DPO loss to optimize alignment.
    \item \textbf{Online Rejection Sampling Fine-tuning (Online RFT)}: Unlike RFT, Online RFT launches the fine-tuning procedure with the SFT model as the initial policy. It subsequently improves this policy by fine-tuning with newly generated outputs from the ongoing, real-time policy model.
    \item \textbf{PPO/GRPO}: In this setting, the initial policy model is derived from the SFT model, and is then optimized by reinforcing it using outputs sampled directly from the current policy model in an online manner.
\end{itemize}

Table \ref{tab:data_policy} provides an overview of the constituent elements of these methods. For an expanded discussion of the derivation steps, consult Appendix \ref{app:analysis-rl}.

\paragraph{Observation about Data Source} We categorize data sources into two types: online sampling and offline sampling. Online sampling refers to data gathered from the interactions generated by the training policy as it learns in real time, whereas offline sampling involves data drawn from the outputs of the preliminary SFT model. Methods such as RFT and DPO utilize offline sampling, whereas Online RFT and GRPO employ online sampling.

Figure \ref{fig:rl-analysis} illustrates that Online RFT achieves substantially better results than RFT on both benchmarks examined. At the outset of training, Online RFT performs at a similar level to RFT; however, as training progresses, Online RFT secures a clear lead, underscoring the benefits of the online learning approach. This observation aligns with expectations: during the early training phase, the behaviors of the actor and the SFT model remain quite similar, with only slight variations in the sampled data. As training advances, larger discrepancies emerge in the samples generated by the actor, making real-time data collection increasingly advantageous.

\paragraph{Observation about Gradient Coefficient} The algorithm utilizes the reward signal to modulate the gradient coefficient, thereby influencing updates to the model parameters. In our experimental setup, the reward function is categorized into 'Rule' and 'Model'. 'Rule' pertains to evaluating the appropriateness of a response by verifying the correctness of the answer, whereas 'Model' involves employing a reward model—trained using data labeled according to the rule-based assessment—to assign scores to each response. Equations \ref{eq:GC-RFT} and \ref{eq:GC-GRPO} illustrate a fundamental distinction between GRPO and Online RFT. Specifically, GRPO adapts its gradient coefficient responsively in accordance with the reward value delivered by the reward model, permitting nuanced reinforcement and penalization depending on the scale of the reward. Conversely, Online RFT lacks such adaptation; it does not impose penalties on incorrect responses and provides a uniform reinforcement strength for all correctly answered responses.

Figure \ref{fig:rl-analysis} illustrates that GRPO achieves better performance than online RFT, underscoring the effectiveness of adjusting the coefficients for positive and negative gradients. Moreover, GRPO+PS outperforms GRPO+OS, suggesting that incorporating granular, step-specific gradient coefficients yields additional advantages. We also examine iterative RL; in our experimental setup, two iterations are performed. Results displayed in Figure \ref{fig:iter-rl} reveal that iterative RL leads to notable performance gains, with the most prominent improvement occurring during the initial iteration.

\subsubsection{Why RL Works?} In our study, we utilize reinforcement learning on a selected portion of instruction tuning datasets, which leads to notable improvements compared to models trained solely with instruction tuning. To elucidate the reasons for this enhancement, we analyze both Pass@K and Maj@K accuracy metrics for the Instruct and RL-based models across two distinct benchmarks. Figure \ref{fig:maj-pass} illustrates that while RL substantially improves the Maj@K metric, it does not yield gains in Pass@K. This pattern suggests that reinforcement learning primarily strengthens the reliability of the output distribution—specifically, \textbf{the observed gains seem to arise from increasing the likelihood of the correct answer being selected among the TopK candidates, rather than from improving core modeling skills.} Along similar lines, \citep{wang2023making} document a \textbf{misalignment problem} regarding reasoning abilities in SFT models, and demonstrate that the reasoning proficiency of SFT models can benefit from various preference alignment methodologies \citep{yuan2023rrhf,song2023preference,wang2023making}.

\subsubsection{How to Achieve More Effective RL?}
\label{sec:effec_rl}
Our results show that RL is highly effective for mathematical reasoning tasks. Furthermore, we introduce a coherent framework that clarifies the connections among various prominent training strategies. In this framework, these strategies are interpreted as forms of either direct or approximated RL approaches. As encapsulated in Equation \ref{eq:objective}, three fundamental elements are involved: Data Source, Algorithm, and Reward Function. We outline several possible avenues for future research related to these components.

\paragraph{Data Source} The foundation of any training approach lies in its data source. Within RL frameworks, we define the data source as consisting of unlabeled questions accompanied by responses generated from the policy model. In this study, our experiments are limited to utilizing questions drawn from the instruction tuning phase, with output generation conducted through simple nucleus sampling. We hypothesize that this limitation may partially account for our RL pipeline's improvements being restricted to Maj@K metrics. For future work, we intend to test our RL pipeline with question prompts outside the original distribution, integrating \textbf{advanced sampling (decoding) strategies}—for instance, leveraging tree-search based techniques \citep{yao2023tree}. Furthermore, the adoption of \textbf{efficient inference techniques} \citep{xia-etal-2023-speculative,leviathan2023fast,kwon2023efficient,xia2024unlocking}, which are instrumental in enhancing the policy model's exploratory capabilities, is anticipated to have a significant impact.

\paragraph{Algorithms} Algorithms utilize both data and the reward signal to compute the gradient coefficient, which is then used for updating model parameters. As articulated in Equation \ref{eq:objective}, present techniques largely \textbf{TRUST} the reward function’s output, adjusting the conditional probability of individual tokens upwards or downwards accordingly. Nonetheless, the reliability of reward signals cannot be fully guaranteed, particularly for tasks with substantial complexity. Notably, datasets such as PRM800K \citep{lightman2023let}, which are the result of meticulous annotation by highly trained experts, still exhibit a mislabeling rate of about 20\%\footnote{\url{https://github.com/openai/prm800k/issues/12\#issuecomment-1728491852}}. Given this limitation, our investigation focuses on reinforcement learning algorithms engineered to tolerate and perform well in the presence of noisy reward feedback. We contend that \textbf{WEAK-TO-STRONG} alignment strategies \citep{burns2023weak} have the potential to radically reshape the landscape of learning algorithms.

\paragraph{Reward Function} The reward function serves as the principal source of learning signals during training. In reinforcement learning, this function typically takes the form of a neural reward model. We identify three key research directions for the advancement of reward models: 1) \textbf{Strategies to bolster the reward model’s generalization capabilities.} To facilitate meaningful improvement in reinforcement learning, it is crucial that the reward model generalizes well, managing out-of-distribution inputs and complex output generations; otherwise, RL could simply regularize the output distribution of LLMs instead of promoting substantive performance gains; 2) \textbf{Methods to capture and represent reward model uncertainty.} Explicit modeling of uncertainty may provide crucial information to bridge the gap between weak reward signals and weak-to-strong training strategies; 3) \textbf{Approaches for constructing robust process-level reward models efficiently} so that detailed and nuanced feedback is available throughout the reasoning process \citep{lightman2023let,wang2023math}.

\section{Conclusion, Limitation, and Future Work}

In this paper, we introduce DeepSeekMath, a model that sets a new standard among open-source models on the MATH competition-level benchmark and narrows the gap with proprietary solutions. DeepSeekMath is based on DeepSeek-Coder-v1.5 7B as its foundation and is continually trained over 500B tokens, notably incorporating 120B mathematical tokens derived from Common Crawl as a core portion of its dataset. Our thorough ablation analysis reveals that web pages are a rich source of quality mathematical content, whereas data from arXiv does not yield as much benefit as initially anticipated. We propose Group Relative Policy Optimization (GRPO), an adaptation of Proximal Policy Optimization (PPO), which substantially enhances mathematical reasoning proficiency while optimizing memory usage. Experimental evidence demonstrates GRPO’s effectiveness, even when DeepSeekMath-Instruct 7B already achieves high benchmark results. Furthermore, we offer a unified framework to interpret a range of approaches and outline key avenues for advancing reinforcement learning methodologies.

While DeepSeekMath~performs notably well on quantitative reasoning benchmarks, it demonstrates comparatively lower proficiency in geometry and theorem-proving tasks than closed-source models. For example, during preliminary testing, the model was unable to address problems involving triangles and ellipses. This limitation potentially points to bias in the data used during pre-training and fine-tuning processes. Moreover, due to constraints related to model size, DeepSeekMath~lags behind GPT-4 in terms of few-shot learning capabilities; whereas GPT-4 benefits from few-shot examples to enhance its accuracy, DeepSeekMath~shows little distinction between zero-shot and few-shot results. Going forward, we plan to enhance our engineered data selection system to assemble a higher-quality pre-training dataset. Additionally, we aim to investigate possible avenues (see Section \ref{sec:effec_rl}) to enable more effective reinforcement learning approaches for large language models.

\end{document}